\begin{ex}
Đường gấp khúc ABC gồm bao nhiêu đoạn thẳng?
\shortans{2}
\loigiai{
Gồm 2 đoạn: AB và BC.
}
\end{ex}

\begin{ex}
Đường gấp khúc MNPQ gồm bao nhiêu đoạn thẳng?
\shortans{3}
\loigiai{
Gồm 3 đoạn: MN, NP, PQ.
}
\end{ex}

\begin{ex}
Hình tứ giác có bao nhiêu cạnh?
\shortans{4}
\loigiai{
Có 4 cạnh.
}
\end{ex}

\begin{ex}
Hình tứ giác có bao nhiêu đỉnh?
\shortans{4}
\loigiai{
Có 4 đỉnh.
}
\end{ex}

\begin{ex}
Hình tam giác có bao nhiêu cạnh?
\shortans{3}
\loigiai{
Có 3 cạnh.
}
\end{ex}

\begin{ex}
Tính độ dài đường gấp khúc gồm 3 đoạn thẳng có độ dài lần lượt là 12 cm, 15 cm, 18 cm (đơn vị cm).
\shortans{45}
\loigiai{
Độ dài: $12 + 15 + 18 = 45$ cm.
}
\end{ex}

\begin{ex}
Tính chu vi hình tam giác có các cạnh dài 8 cm, 10 cm, 12 cm (đơn vị cm).
\shortans{30}
\loigiai{
Chu vi: $8 + 10 + 12 = 30$ cm.
}
\end{ex}

\begin{ex}
Tính chu vi hình tứ giác có các cạnh dài 5 cm, 7 cm, 9 cm, 11 cm (đơn vị cm).
\shortans{32}
\loigiai{
Chu vi: $5 + 7 + 9 + 11 = 32$ cm.
}
\end{ex}

\begin{ex}
Một hình vuông có cạnh dài 8 cm. Chu vi hình vuông là bao nhiêu xăng-ti-mét?
\shortans{32}
\loigiai{
Chu vi: $8 \times 4 = 32$ cm.
}
\end{ex}

\begin{ex}
Một hình chữ nhật có chiều dài 10 cm, chiều rộng 6 cm. Nửa chu vi là bao nhiêu xăng-ti-mét?
\shortans{16}
\loigiai{
Nửa chu vi: $10 + 6 = 16$ cm.
}
\end{ex}

\begin{ex}
Chu vi hình chữ nhật đó là bao nhiêu xăng-ti-mét?
\shortans{32}
\loigiai{
Chu vi: $16 \times 2 = 32$ cm.
}
\end{ex}

\begin{ex}
Cho 3 điểm thẳng hàng A, B, C (B nằm giữa A và C). Biết AB dài 7 cm, BC dài 8 cm. Độ dài đoạn AC là bao nhiêu xăng-ti-mét?
\shortans{15}
\loigiai{
Độ dài AC: $7 + 8 = 15$ cm.
}
\end{ex}

\begin{ex}
Đoạn thẳng MN dài 20 cm, điểm P nằm giữa M và N sao cho MP dài 9 cm. Đoạn thẳng PN dài bao nhiêu xăng-ti-mét?
\shortans{11}
\loigiai{
Đoạn PN dài: $20 - 9 = 11$ cm.
}
\end{ex}

\begin{ex}
Cần ít nhất bao nhiêu que tính để xếp được 1 hình tứ giác?
\shortans{4}
\loigiai{
Cần 4 que tính.
}
\end{ex}

\begin{ex}
Cần ít nhất bao nhiêu que tính để xếp được 1 hình tam giác?
\shortans{3}
\loigiai{
Cần 3 que tính.
}
\end{ex}

\begin{ex}
Dùng 9 que tính bằng nhau xếp được nhiều nhất bao nhiêu hình tam giác rời nhau?
\shortans{3}
\loigiai{
Xếp được: $9 : 3 = 3$ hình tam giác.
}
\end{ex}

\begin{ex}
Dùng 12 que tính bằng nhau xếp được nhiều nhất bao nhiêu hình vuông rời nhau?
\shortans{3}
\loigiai{
Xếp được: $12 : 4 = 3$ hình vuông.
}
\end{ex}

\begin{ex}
Một đường gấp khúc dài 36 cm gồm 3 đoạn thẳng bằng nhau. Mỗi đoạn dài bao nhiêu xăng-ti-mét?
\shortans{12}
\loigiai{
Mỗi đoạn dài: $36 : 3 = 12$ cm.
}
\end{ex}

\begin{ex}
Một hình tứ giác có chu vi 48 cm, bốn cạnh có độ dài bằng nhau. Mỗi cạnh dài bao nhiêu xăng-ti-mét?
\shortans{12}
\loigiai{
Mỗi cạnh dài: $48 : 4 = 12$ cm.
}
\end{ex}

\begin{ex}
Cắt một góc của hình tứ giác bằng một nhát kéo thẳng. Hình mới nhận được có bao nhiêu đỉnh?
\shortans{5}
\loigiai{
Hình mới có 5 đỉnh.
}
\end{ex}

\begin{ex}
Gấp đôi một hình vuông theo đường chéo ta được hình gì? Hình tam giác mới có bao nhiêu cạnh?
\shortans{3}
\loigiai{
Hình tam giác có 3 cạnh.
}
\end{ex}

\begin{ex}
Một con sâu bò trên đường gấp khúc gồm 4 đoạn, mỗi đoạn dài 6 cm. Con sâu bò được quãng đường bao nhiêu xăng-ti-mét?
\shortans{24}
\loigiai{
Quãng đường: $6 \times 4 = 24$ cm.
}
\end{ex}

\begin{ex}
Hình chữ nhật có bao nhiêu góc vuông?
\shortans{4}
\loigiai{
Có 4 góc vuông.
}
\end{ex}

\begin{ex}
Hình vuông có bao nhiêu góc vuông?
\shortans{4}
\loigiai{
Có 4 góc vuông.
}
\end{ex}

\begin{ex}
Tổng số cạnh của 2 hình tam giác rời nhau là bao nhiêu?
\shortans{6}
\loigiai{
Tổng số cạnh: $3 + 3 = 6$ cạnh.
}
\end{ex}
